\section{Clinical Background}
\label{sec:clinical}

\subsection{Cervical anatomy on sagittal T2 MRI}
The cervical spine comprises seven vertebrae (C1--C7) and the intervening discs (C2--C3 through
C7--T1). On a mid-sagittal T2 image, the bright cerebrospinal fluid (CSF) outlines the thecal sac, with
the spinal cord visible as an intermediate-signal band within it. We measure C3--C7 and the discs between
them; C1 (atlas) and C2 (axis/odontoid) are excluded because their morphology is structurally unique and
the standard vertebral-body measurements do not apply.

\subsection{The measurements and what they mean clinically}
\paragraph{Vertebral body (Group 1).} Anterior, middle, and posterior heights ($H_a$, $H_m$, $H_p$);
the compression ratio \hahp{}; antero-posterior (AP) width; tilt; and spondylolisthesis (slip). A
reduced \hahp{} indicates anterior wedging (compression/deformity); a Genant-style grade describes the
degree~\cite{genant1993}. Spondylolisthesis is forward/backward slip of one vertebra on another.

\paragraph{Intervertebral disc (Group 2).} Disc signal-intensity height, the disc-height index (DHI),
posterior disc bulge, and a signal-degeneration grade. Disc-height loss and signal loss are hallmarks of
degeneration; a posterior bulge or herniation can compress the cord or nerve roots. Cervical disc signal
is graded by the modified (Miyazaki) scheme~\cite{miyazaki2008}, not the lumbar Pfirrmann scheme.

\paragraph{Canal and cord (Group 3).} The functional (soft-tissue) canal / dural-sac AP minimum, the
spinal cord AP diameter, the space available for the cord (SAC $=$ canal $-$ cord), and the
Torg--Pavlov ratio. Canal narrowing and a small SAC are the structural substrate of myelopathy; the
cord may flatten under compression. Normative cervical canal, SAC, and vertebral-width values are
available from population studies~\cite{nell2019}, and cord cross-section can be compared to a healthy
normative database~\cite{valosek2024}.

\paragraph{Sagittal alignment (Group 4).} The C2--C7 (and C3--C7) Cobb angle, segmental angles, and a
lordosis classification. The cervical spine is normally lordotic; loss of lordosis or kyphosis is
associated with degeneration and with myelopathy outcomes. Asymptomatic C2--C7 lordosis is roughly
\SIrange{10}{35}{\degree}~\cite{nassj2025}.

\paragraph{Signal and shape screens (Group 5).} A vertebral-body compression/deformity screen and a cord
T2-hyperintensity (myelomalacia) screen. Intramedullary T2 hyperintensity is a marker of myelopathy.

\subsection{Two clinical syndromes the cohort represents}
Cervical spondylosis (degenerative ``wear-and-tear'') produces two principal syndromes that our
unhealthy cohort is labelled by:
\begin{itemize}[leftmargin=*]
  \item \textbf{Cervical Spondylotic Myelopathy (CSM)} --- spinal-cord dysfunction from canal stenosis
  and cord compression. Structurally this is a Group 3 (canal/SAC/cord) and Group 5.1 (cord signal)
  disease.
  \item \textbf{Cervical Spondylotic Radiculopathy (CSR)} --- nerve-root compression, usually from disc
  herniation/osteophyte. Structurally this is a Group 2 (disc) disease.
\end{itemize}
Crucially, \emph{spondylosis is not the same as vertebral compression fracture}. Compression fractures
arise from osteoporosis or trauma, not degeneration, so a spondylosis cohort is expected to have normal
vertebral-body heights. This distinction (spondyl\emph{osis} vs.\ spondyl\emph{olisthesis} vs.\
compression) is central to interpreting our validation: it explains why the Group 1 compression screen
should be, and is, silent on this cohort (Section~\ref{sec:results}).

\subsection{The no-diagnosis constraint}
Quantitative measurements are necessary but not sufficient for diagnosis: an abnormal value can occur in
asymptomatic individuals (for example, a large fraction of asymptomatic adults have a measurable disc
bulge or a low Torg ratio), and clinical correlation is always required. The system therefore reports
measurements and threshold crossings as findings ``flagged for physician review,'' never as diagnoses ---
a hard rule enforced throughout the codebase and the output wording.
